\begin{center}
  \large\textbf{ABSTRACT}
\end{center}

\addcontentsline{toc}{chapter}{ABSTRACT}

\vspace{2ex}

\begingroup
% Menghilangkan padding
\setlength{\tabcolsep}{0pt}

\noindent
\begin{tabularx}{\textwidth}{l >{\centering}m{3em} X}
  \emph{Name}     & : & \name{}         \\

  \emph{Title}    & : & \engtatitle{}   \\

  \emph{Advisors} & : & 1. \advisor{}   \\
                  &   & 2. \coadvisor{} \\
\end{tabularx}
\endgroup

% Ubah paragraf berikut dengan abstrak dari tugas akhir dalam Bahasa Inggris
\emph{Electronic ticketing systems provide convenience in ticket purchasing and management while reducing the risk of ticket loss. However, centralization remains a challenge, including vulnerabilities to cyberattacks and dependence on third parties. Blockchain technology offers a solution by providing a decentralized, transparent, and secure transaction ledger. It also supports the development of digital assets such as NFTs, opening new opportunities for digital transactions and collections. To overcome the limitations of cost and scalability on Ethereum’s main network, Polygon zkEVM emerges as a Layer-2 solution utilizing zero-knowledge rollups, enabling more efficient, faster, and cheaper transactions without compromising security. Although challenges such as intellectual property rights and cross-chain interoperability remain, the adoption of blockchain and Polygon zkEVM demonstrates great potential in building an inclusive and efficient digital asset ecosystem.}

% Ubah kata-kata berikut dengan kata kunci dari tugas akhir dalam Bahasa Inggris
\emph{Keywords}: \emph{Blockchain}, \emph{Electronic Ticketing}, \emph{NFT}, \emph{Polygon
  zkEVM}, \emph{Decentralization}.
